\section{Kernel-Level Profile and Alternative Execution Models}
\label{app:profile}

This appendix opens the iteration with the vendor profilers and measures the
execution models a systems reader would propose in place of graph replay: more
iterations per graph, a generic tensor compiler, more CPU cores, and graph capture
of the prior GPU baseline.

\paragraph{What one iteration launches.}
\cref{tab:profile} traces 100 steady-state iterations with Nsight Systems inside an
NVTX window, in eager and graph mode, and replays two eager iterations under Nsight
Compute. The Aten operation count of \cref{tab:suite} is close to, but not the same
as, the kernel count: on the turn subgame 96 operations become 87 kernels (some
operations dispatch no kernel, a few dispatch two), and the graph path replays the
same kernels plus the device-resident counter update. Eager execution submits 87
host launches per iteration and keeps a kernel resident on the GPU for under half of
the window even with the tracer's own launch overhead inflating the iteration. Graph
replay issues one graph launch per iteration; the two further host launches per
iteration in \cref{tab:profile} are the fill kernels with which PyTorch's
\texttt{CUDAGraph.replay()} refreshes its random-number state on every replay, a
framework fixture the solver never reads, and the three extra kernels inside the
graph are those two fills plus the device-side iteration-counter increment. GPU
residency rises above 90\%. The remaining time is kernel execution: 87 kernels averaging under
5\,$\mu$s each, dominated by the gather kernels of the forward and backward passes
(38 per iteration, 58\% of GPU time), then the scatter and \texttt{index\_add}
kernels of the regret and strategy accumulation. Half of the kernels launch fewer
blocks than the A100 has SMs, so the device is latency bound on small grids rather
than throughput bound: the counters put DRAM traffic at 129\,MB per iteration, which
at the unprofiled 0.397\,ms is about 17\% of peak bandwidth averaged over the
iteration and 9\% time-weighted inside kernels, and SM throughput at 5\% of peak.
The bytes-moved model behind \cref{fig:scaling} counts only the solver's own
arrays and therefore reports a lower bound on this traffic.

\paragraph{More iterations per graph.}
Because the iteration counter advances on the device, any number of iterations can
be recorded into one graph (\cref{sec:cudagraph}). \cref{tab:altexec} records 1, 10,
and 100 iterations per graph: the per-iteration time is unchanged within
measurement noise on all three games. A graph launch costs a few microseconds
against a 0.304--0.557\,ms body, so the single-iteration graph already captures the
whole launch-amortization benefit, and the remaining cost is kernel execution.

\paragraph{A generic compiler on the same dataflow.}
\cref{tab:altexec} also re-measures the ladder under PyTorch 2.5.1+cu121 in one
process and wraps the compiled iteration body in \texttt{torch.compile} (Inductor),
driven with the same device-resident counter so the trace contains no per-iteration
Python scalars. The three ladder stages keep their PyTorch 1.13.1 ordering and magnitudes
(19$\times$ and 1.9$\times$ against 19$\times$ and 1.6$\times$), which settles the
framework-version question for \cref{fig:ladder}. On the turn subgame Inductor's
\texttt{default} mode runs 11\% faster than CUDA graph replay of the uncompiled body
(0.357 against 0.397\,ms); on the launch-bound Leduc tree it stays 1.4$\times$
slower, and \texttt{reduce-overhead} mode (which itself captures CUDA graphs) is no
faster in either case. The gather and scatter kernels that carry the iteration are
not fused into fewer launches by any tested Inductor configuration, so a generic compiler recovers little beyond the launch
amortization that graph replay already provides. It does so at a price the graph
path does not pay: the first call compiles for seconds (4.025--4.747\,s on the turn
subgame), and the compiled floats differ from the eager iterates because Inductor
reorders reductions, whereas graph replay runs the identical kernels in the
identical order. The game-level compilation of
\cref{sec:method} is what turns the tree walk into this flat dataflow in the first
place; the generic compiler can only consume the result.

\paragraph{Graph capture of the prior GPU baseline.}
Kim (2026)'s iteration cannot be captured. Its CuPy sparse products call cuSPARSE,
which CuPy refuses to execute under stream capture, and its per-level row slicing,
CSR index assignment, and boolean gathers each perform a host-to-device transfer per
iteration that capture forbids. We verified both on a Kuhn-sized instance:
\texttt{cupy.cuda.Stream.begin\_capture} raises on the first sparse product. The
representation is what makes the iteration capturable.

\paragraph{CPU thread scaling.}
\cref{fig:cpuscaling} (\cref{sec:results}) sweeps the compiled CPU solver from 1 to 28 threads pinned to
whole physical cores of one socket on the four largest games. Throughput saturates
well below the core count: the 8-thread arm of \cref{tab:matrix} sits within a
small factor of the 28-thread point on every game, because each depth block is a
gather or scatter over a few hundred thousand elements that PyTorch's intra-operator
pool cannot split efficiently. The 8-thread arm is therefore a representative CPU
ceiling for this dataflow.


\paragraph{Hand-fused kernels.}
The remaining execution model is to replace the generic gather and scatter kernels
with kernels written for this dataflow. \Cref{tab:kernels} measures a Triton
implementation that computes one depth level of the forward or backward pass per
launch and fuses regret matching, the regret update and the average-strategy update
into one kernel each, so an iteration issues $2D + 5$ launches against the
$8D + 32$ Aten operations of the generic path, the same count \cref{tab:suite} reports
(recounted under the PyTorch 2.5 dispatcher, it is unchanged); it drives the solver's own buffers and
index arrays and agrees with the generic path to $10^{-9}$ in float64. All four arms
run in one process on the PyTorch 2.5 stack that ships Triton, both with and without
graph capture, so the table separates the gain of fusion from the gain of launch
batching and places the compiled tensor path with graph replay on that frontier. The
compiled path keeps the property that motivates it: the same code runs every update
rule, every game, and the CPU, from one representation and without a kernel per
dataflow.

\input{figures/TABLE_T11_profile.tex}
\input{figures/TABLE_T12_altexec.tex}
\input{figures/TABLE_T15_kernels.tex}
